%===============================================================================
\section{Related Work}\label{sec:related}
%===============================================================================

\paragraph{Inference and aggregation in data release}
That individually innocuous data items can jointly reveal protected information is an old
observation in database security, where it is known as the inference and aggregation
problem~\cite{lunt1989aggregation,farkas2002inference}; the privacy literature supplies the
best-known illustration, three demographic attributes that jointly single out most
individuals~\cite{sweeney2000simple}. In power systems, public market data have been shown to
expose non-public physical information, for example grid topology from locational marginal
prices~\cite{kekatos2014topology}. Recent work studies reconstruction through attribute correlations in multi-attribute releases~\cite{jiang2025correlation}. This is closely related to the motivation for measuring conditional contributions, while our audit derives a field profile from the inferability of admissible sets. Privacy mechanisms for metering data~\cite{mcdaniel2009smartgrid,giaconi2018smartmeter}
and formal models such as differential privacy~\cite{dwork2014dp} or
$k$-anonymity~\cite{sweeney2002kanonymity} bound what can be learned about an individual
record after a release mechanism is applied. They answer a different question from the one
faced by a market operator, who must decide which aggregate operational \emph{fields} may be
published at all.

\paragraph{Data sensitivity assessment under inference leakage}
Data classification and grading is performed field by field in practice, and the Chinese
national standard explicitly asks that aggregation and processing effects be
considered~\cite{gbt43697}, without prescribing how. The closest work to ours is IGNN~\cite{hu2026ignn},
which introduced data inference leakage into sensitivity assessment for cyber-physical power
grids: fields are nodes of an inference graph whose edges encode physical or statistical
inference relations, and sensitivity is propagated along these edges by a graph attention
network from an intrinsic (expert-assigned) score. This established that the risk of a field
depends on its inference relations to sensitive data. Two questions remain open in this line
of work. First, the risk score is calibrated against expert labels, while its ultimate
validation is an attack experiment in which a model is trained on a group of fields; the
attack outcome itself is not the definition. Second, an edge connects two fields, whereas
inference frequently requires several fields jointly, with none of them informative alone.
We take the attack outcome on a field \emph{set} as the definition of risk and derive the
field-level quantity from it.

\paragraph{Feature importance}
Turning set-level values into per-feature scores is the subject of feature attribution.
Shapley-based measures~\cite{shapley1953value,lundberg2017shap,covert2020sage,williamson2020spvim}
average a feature's marginal contribution over all contexts, which is the right answer for
fair allocation but dilutes the credit of redundant features~\cite{kumar2020shapley,catav2021mci}.
Leave-one-covariate-out (LOCO)~\cite{lei2018loco} and permutation importance~\cite{breiman2001rf,fisher2019mcr,hooker2021unrestricted}
evaluate a single context, the full feature set, so a feature with a near-substitute receives
almost no importance. The Marginal Contribution Index (MCI)~\cite{catav2021mci} takes the
\emph{maximum} marginal contribution over contexts, is characterized axiomatically, and is
immune to dilution. Security reasoning favours the worst case for the same reason
quantitative information flow replaced Shannon entropy by min-entropy and gain
functions~\cite{smith2009qif,alvim2012gleakage}: an average over contexts is not a guarantee.
Our field-level risk is a budget-constrained MCI evaluated relative to already-public data.

\paragraph{Fast estimation of many refitted models}
Computing the performance of a model refitted on many feature subsets is expensive, which has
motivated lazy~\cite{gao2022lazyvi} and warm-started~\cite{sun2025warmstart} estimators of
variable importance, surrogate models trained under random masks for removal-based
explanations~\cite{covert2021explaining,jethani2022fastshap,covert2023vitshapley}, conditional generative models trained under
random masks~\cite{ivanov2019vaeac,li2020acflow,yoon2020vime,he2022mae}, and tabular foundation models that predict from any
set of columns in context~\cite{hollmann2025tabpfn}. The last amortize learning across data sets: each query passes the
training rows of the requested columns through a large pre-trained transformer. Our estimator amortizes across the field
sets of one data set instead, so that the per-set work reduces to a linear solve on shared features.
Closed-form ridge heads on frozen features are the standard fast-adaptation device in
meta-learning~\cite{bertinetto2019r2d2,lee2019metaoptnet}, and the benefit of sharing a
representation across tasks is well understood~\cite{caruana1997multitask,baxter2000inductive,maurer2016benefit}.
We combine these ingredients into an estimator of set-level inferability and, more
importantly, use it inside a scan--certify loop so that decisions rest on retrained
attackers rather than on the estimate.

%===============================================================================

\paragraph{Efficiency evaluated against the security task}
Practical security methods must account for the cost of producing useful evidence. For example, LPASS uses linear probes to estimate the performance of compressed models for vulnerability detection and evaluates both security-task performance and computational savings~\cite{ibanez2025lpass}. Our application has a different workload: thousands of related field subsets within one data set. The shared estimator is evaluated against retrained inferability and downstream background recovery, rather than against prediction accuracy alone.
